\documentclass[10pt]{article}
\usepackage[T1]{fontenc}
\usepackage[utf8]{inputenc}
\usepackage[french]{babel}
\usepackage{amsmath,amssymb,mathrsfs,mathpazo}
\usepackage{geometry}
\geometry{a4paper,margin=2cm}
\usepackage{graphicx,array,tabularx,tabularray,colortbl,multirow,multicol,booktabs,enumitem,pifont,ulem,fancybox,diagbox,cancel}
\usepackage{pstricks,pstricks-add,pst-plot,pst-tree,pst-node,pst-text,pst-func,pst-math,pst-bspline,pst-3dplot,pst-eucl,pst-solides3d}
\let\pscalculate\undefined
\usepackage{pst-bezier}
\usepackage{tikz,pgfplots,tkz-tab}
\usetikzlibrary{arrows,arrows.meta,patterns,calc,positioning,trees,decorations.markings,shapes.geometric,shapes.misc,shapes.arrows,intersections,decorations.pathmorphing,decorations.pathreplacing,backgrounds,shadows,babel}
\pgfplotsset{compat=1.18}
\usepgfplotslibrary{fillbetween}
\usepackage[np]{numprint}
\usepackage{listings,makecell,pas-tableur,twoopt}
\let\barre\undefined
\usepackage[pstricks]{tablvar}
\definecolor{darkgreen}{RGB}{0,100,0}
\providecommand{\red}{\color{red}}
\providecommand{\blue}{\color{blue}}
\providecommand{\green}{\color{green}}
\providecommand{\black}{\color{black}}
\providecommand{\white}{\color{white}}
\providecommand{\gray}{\color{gray}}
\providecommand{\cyan}{\color{cyan}}
\providecommand{\magenta}{\color{magenta}}
\providecommand{\yellow}{\color{yellow}}
\providecommand{\violet}{\color{violet}}
\providecommand{\orange}{\color{orange}}
\providecommand{\euro}{\texteuro}
\providecommand{\EUR}{\texteuro}
\providecommand{\ieme}{\textsuperscript{e}}
\providecommand{\degres}{\ensuremath{^\circ}}
\providecommand{\degree}{\ensuremath{^\circ}}
\providecommand{\vv}[1]{\overrightarrow{#1}}
\providecommand{\Coord}[1]{#1}
\providecommand{\decosix}{}
\providecommand{\starredbullet}{\ensuremath{\bullet}}
\pagestyle{empty}
\setlength{\parindent}{0pt}
\usepackage{adjustbox}\DeclareUnicodeCharacter{2236}{:}\DeclareUnicodeCharacter{2190}{\ensuremath{\leftarrow}}
\providecommand{\R}{\mathbb{R}}
\providecommand{\N}{\mathbb{N}}
\providecommand{\Z}{\mathbb{Z}}
\providecommand{\Q}{\mathbb{Q}}
\providecommand{\D}{\mathbb{D}}
\providecommand{\C}{\mathbb{C}}
\providecommand{\e}{{\mathrm{e}}}
\providecommand{\vect}[1]{\overrightarrow{#1}}
\providecommand{\vectt}[1]{\overrightarrow{\mathrm{#1}}}
\providecommand{\vvt}[1]{\overrightarrow{\mathrm{#1}}}
\providecommand{\np}[2][]{{#2~#1}}
\providecommand{\vv}[1]{{\overrightarrow{#1}}}
\renewcommand{\d}{{\mathrm{d}}}
\providecommand{\degres}{{^\circ}}
\providecommand{\degre}{{^\circ}}
\providecommand{\blue}{\color{blue}}
\providecommand{\red}{\color{red}}
\providecommand{\enskip}{\quad}
\providecommand{\llbracket}{\lbrack\!\lbrack}
\providecommand{\rrbracket}{\rbrack\!\rbrack}
\providecommand{\euro}{{\text{€}}}
\providecommand{\no}{{\text{n°}}}
\providecommand{\interval}[2]{{[#1;#2]}}
\providecommand{\intervFO}[2]{{[#1;#2[}}
\providecommand{\dotfill}{\ldots}

\providecommand{\ds}{\displaystyle}
\providecommand{\textding}[1]{#1}
\providecommand{\Oij}{\(\mathrm{O};\vec{\imath},\vec{\jmath}\)}
\providecommand{\Oijk}{\(\mathrm{O};\vec{\imath},\vec{\jmath},\vec{k}\)}
\providecommand{\Ouv}{\(\mathrm{O};\vec{u},\vec{v}\)}
\providecommand{\intervFF}[2]{\ensuremath{[#1;#2]}}
\renewcommand{\euro}{\ifmmode\text{\texteuro}\else\texteuro\fi}
\DeclareUnicodeCharacter{00B0}{\ifmmode^{\circ}\else\textdegree\fi}
\DeclareUnicodeCharacter{20AC}{\ifmmode\text{\texteuro}\else\texteuro\fi}
% Énoncés : Éducation nationale. Transcription et corrigés : APMEP, Denis Vergès et contributeurs. Découpage et classement : CoopMaths, Sébastien Lozano, Éric Elter et contributeurs. Adaptation de la présentation : Maths Explorer.
% CC BY-SA : https://coopmaths.fr/www/about/
% https://coopmaths.fr/alea/static/bac/2025/tex/bac_2025_09_sujet1_metropole_3.tex
\begin{document}
\noindent\textit{APMEP / CoopMaths. CC BY-SA. Présentation adaptée par Maths Explorer.}\par\medskip
\medskip

Le but de cet exercice est d'étudier les convergences de deux suites vers une même limite.

\medskip

\textbf{Partie A}

\medskip

On considère la fonction $f$ définie sur $[2~;~+\infty[$ par 

\[f(x) = \sqrt{3x - 2}.\]

\medskip

\begin{enumerate}
\item Justifier les éléments du tableau de variations ci-dessous :











\[\begin{tablvar}[stretch=1.2,intervalwidth=8em]{1}
\hline
x & 2 && +\infty\\
\hline
\variations{\mil{f(x)} & \bas{2} && \haut{+\infty}}
\hline
\end{tablvar}\]

On admet que la suite $\left(u_n\right)$ vérifiant $u_0 = 6$ et, pour tout $n$ entier naturel, $u_{n+1} = f\left(u_n\right)$ est bien définie.
\item 
	\begin{enumerate}
		\item Démontrer par récurrence que, pour tout $n$ entier naturel : $2 \leqslant u_{n+1} \leqslant u_n \leqslant 6$.
		\item En déduire que la suite $\left(u_n\right)$ converge.
	\end{enumerate}
\item On appelle $\ell$ la limite de $\left(u_n\right)$.

On admet qu'elle est solution de l'équation $f(x) = x$. Déterminer la valeur de $\ell$.

\item On considère la fonction \texttt{rang} écrite ci-dessous en langage Python.

On rappelle que sqrt($x$) renvoie la racine carrée du nombre $x$.

\begin{center}
\begin{ttfamily}
\begin{tabular}{>{\cellcolor{lightgray!50}}c l}
1&from math import *\\
2&\\
3&def rang(a) :\\
4&\quad  u = 6\\
5&\quad n=0\\
6& \quad while u > = a :\\
7&\qquad u = sqrt(3*u - 2)\\
8&\qquad n = n+1\\
9& \quad return n\\
\end{tabular}
\end{ttfamily}
\end{center}

	\begin{enumerate}
		\item Pourquoi peut-on affirmer que \texttt{rang}(2.000001) renvoie une valeur ?
		\item Pour quelles valeurs du paramètre $a$ l'instruction \texttt{rang}($a$) renvoie-t-elle un résultat ?
	\end{enumerate}
\end{enumerate}

\bigskip

\textbf{Partie B}

\medskip

On admet que la suite $\left(v_n\right)$ vérifiant $v_0 = 6$ et, pour tout $n$, entier naturel, $v_{n+1} = 3 - \dfrac{2}{v_n}$ est bien définie.

\medskip

\begin{enumerate}
\item Calculer $v_1$.
\item Pour tout $n$ entier naturel, on admet que $v_n \ne 2$ et on pose :

\[w_n = \dfrac{v_n - 1}{v_n - 2}\]

	\begin{enumerate}
		\item Démontrer que la suite $\left(w_n\right)$ est géométrique de raison 2 et préciser son premier terme $w_0$.
		\item On admet que, pour tout $n$ entier naturel,

\[w_n - 1 = \dfrac{1}{v_n - 2}.\]

En déduire que, pour tout $n$ entier naturel,

\[v_n = 2 + \dfrac{1}{1,25 \times 2^n - 1}\]

		\item Calculer la limite de $\left(v_n\right)$.
	\end{enumerate}
\item Déterminer le plus petit entier naturel $n$ pour lequel $v_n < 2,01$ en résolvant
l'inéquation.
\end{enumerate}

\bigskip

\textbf{Partie C}

\medskip

À l'aide des parties précédentes, déterminer le plus petit entier $N$ tel que pour tout $n \geqslant N$, les termes $v_n$ et $u_n$ appartiennent à l'intervalle ]1,99~;~2,01[.

\bigskip

\end{document}
